%% LyX 2.0.8.1 created this file.  For more info, see http://www.lyx.org/.
%% Do not edit unless you really know what you are doing.
\documentclass[english,aip, reprint]{revtex4-1}
\usepackage{charter}
\usepackage[T1]{fontenc}
\usepackage[latin9]{inputenc}
\setcounter{secnumdepth}{3}
\usepackage{color}
\usepackage{babel}
\usepackage{amsbsy}
\usepackage{graphicx}
\usepackage{esint}
\usepackage[unicode=true,pdfusetitle,
 bookmarks=true,bookmarksnumbered=false,bookmarksopen=false,
 breaklinks=false,pdfborder={0 0 0},backref=false,colorlinks=true]
 {hyperref}
\hypersetup{
 citecolor=blue, urlcolor=blue, linkcolor=blue}

\makeatletter

%%%%%%%%%%%%%%%%%%%%%%%%%%%%%% LyX specific LaTeX commands.
%% A simple dot to overcome graphicx limitations
\newcommand{\lyxdot}{.}


\@ifundefined{showcaptionsetup}{}{%
 \PassOptionsToPackage{caption=false}{subfig}}
\usepackage{subfig}
\makeatother

\begin{document}

\title{Non-linear absorption of intense bi-color laser light by indirect
band gap materials resolved by time-dependent density functional theory}


\author{Thibault J.-Y. Derrien}


\affiliation{Max Planck Institute for Structure and Dynamics of Matter (MPSD),
Hamburg, Germany}


\affiliation{HiLASE Centre, Institute of Physics, Academy of Science of the Czech
Republic, Za Radnic� 828/5, 25241 Doln� B\v{r}e�any, Czech Republic}


\author{Heiko Appel}


\affiliation{Max Planck Institute for Structure and Dynamics of Matter (MPSD),
Hamburg, Germany}


\author{Nicolas Tancogne-Dejean}


\affiliation{Max Planck Institute for Structure and Dynamics of Matter (MPSD),
Hamburg, Germany}


\author{Angel Rubio}


\affiliation{Max Planck Institute for Structure and Dynamics of Matter (MPSD),
Hamburg, Germany}


\author{Nadezhda M. Bulgakova}


\affiliation{HiLASE Centre, Institute of Physics, Academy of Science of the Czech
Republic, Za Radnic� 828/5, 25241 Doln� B\v{r}e�any, Czech Republic}
\begin{abstract}
Time-Dependent Density Functional Theory (TD-DFT) is used to generate
the transition probability of electrons from valence band (VB) to
conduction band (CB) in indirect band-gap semiconductor Si. Two orientations
of crystalline lattice are compared. Typical pulses used in laser
applications are built using a space of 7 tuning parameters, including
wavelength, pulse duration, phases inside the envelope, and the polarization
of laser fields. Purpose of this document is for now to analyze and
keep a trace of first quantum simulations. This is aimed to prepare
a paper on the non-linear absorption of light by Si atoms exposed
to single-color and bi-color laser excitations. 
\end{abstract}
\maketitle
\tableofcontents{}


\section{Introduction}

It was experimentally demonstrated that a dense electron-hole plasma
is generated upon femtosecond laser irradiation \cite{Sokolowski-Tinten2000}.
 Such transient excitation leads to the change of refractive index
as a function of time and laser energy and can lead to excitation
of several processes such as triggering chemical reaction, non-thermal
phase transitions, and surface plasmon polariton excitation \cite{Derrien2013,Sato2014a}.
Making use of multi-photonic absorption also enabled outstanding applications
such as 2-photon polymerization {[}REF{]}, and multiphoton microscopy
{[}REF{]}. 

First theoretical description of non-linear absorption in matter was
successful but are nowadays not sufficient to address the needs for
describing laser-induced band gap energy levels, non-parabolic band
structure, or the wide spectrum of ultrashort laser pulses. Nowadays,
strong hypothesis such as monochromaticity of waves, parabolic band
gap approximation, and direct interband transitions \cite{Keldysh1965,Garcia2006}
could not be easily removed. Therefore, calculating the non-linear
excitation transition rates in matter remain a puzzling problem. 

In addition, making use of several laser pulses together lead to give
more control over laser ablation , nanostructure formation \cite{Derrien2013a,Hoehm2013a,Hoehm2015,Hoehm2015b},
high harmonic generation and {[}see results from Les Houches workshop{]}. 

In view of bridging large scale simulations \cite{Bulgakova2013,Bulgakova2014,Buschlinger2014}
with first principle calculations \cite{Flick2015a,Ruggenthaler2014},
this paper propose a quantum simulation of the non-linear absorption
in indirect band gap material with the example of silicon, based on
the method of TD-DFT. 


\section{Calculation method}

In the quantum modeling world, the linear absorption in bulk is typically
resolved using the meta-Generalized Gradient Approximation (MGGA)
qualitatively well reproducing the experimentally measured absorption
spectrum for Si and Ge \cite{Nazarov2011}. A full picture of the
photo-emission from atoms and molecules also came to be available
\cite{DeGiovannini2012}. 

Hence, it appears that linear absorption can be calculated using methods
like MGGA, or Kohn-Sham approach with exchange-correlation derivated
from Bethe-Salpeter equation. 

Difficulty appears to be linked with the description of excitons bounded
on large computational cell, and with non-linear excitation (intense
fields). So far, the case of gas phase ionization was well studied
(so called ``above threshold ionization'') {[}\cite{Keldysh1965},
\cite{Chin1988} \cite{Ruggenthaler2008,Popruzhenko2008a}{]}. 

Today the difficulty remains to calculate excitation rates in correlated
matter upon intense laser light. Such calculation started for small
metallic systems \cite{Otobe2004,Shiratori2005,Otobe2008,Otobe2009,Tong2015}.
In addition, the effect of intense constant field got recently studied
\cite{Otobe2016}. 

In this paper, we will show results of quantum simulations made on
Si, from the most basic case to the most advanced. This will support
the development of my work in MPSD. We will select the interesting
ones before publication. 


\subsection{Preparation of the laser pulses}


\subsubsection{Lasers in quantum simulations}

As the scale of the sample is of few atomic layers, the local electric
field of the laser can be approximated by its simple temporal variation.
Simulations are resolved in a space-time-dependent approach, hence
the dispersion of the laser pulse is natively taken into account.
For short pulse durations, the spectrum of the pulses becomes wider. 

If using QED-DFT \cite{Flick2015a}, some modulations of the electric
field response may be observed. 


\subsubsection{Temporal evolution of bi-color pulses}


\paragraph{Experimental techniques}

In the experiments, two configurations are possible. 
\begin{description}
\item [{(i)}] We can make use of one laser source, and place a BBO crystal
in the optical path, then focusing both pulses using the same lens,
in the most ideal case. In this case, the two pulses will have blocked
relative phase $CEP_{2}-CEP_{1}$. 
\item [{(ii)}] Second possibility is to use two different lasers, which
may then have two different phases $CEP_{1}$ and $CEP_{2}$, even
maybe time-varying relative phases. 
\end{description}
We now draw a general model which takes into account all possible
configurations. At the atomic scale, the spatial variation of the
pulse envelope can be neglected. Hence, the temporal envelope of the
pulse can be defined by the general bi-color pulse expression
\[
\boldsymbol{E}(t)=\boldsymbol{E}_{\lambda_{1}}(t)+\boldsymbol{E}_{\lambda_{2}}(t,\Delta t).
\]
Each pulse can be developed in
\[
\boldsymbol{E}_{\lambda_{1}}=\boldsymbol{E}_{0}^{(1)}\sin\left(\omega_{1}t+CEP_{1}\right)\cdot\exp\left[-\frac{1}{2}\left(\frac{t}{\sigma\left(\tau_{1}\right)}\right)^{2}\right]
\]
\[
\boldsymbol{E}_{\lambda_{2}}=\boldsymbol{E}_{0}^{(2)}\sin\left(\omega_{2}t+CEP_{2}\right)\cdot\exp\left[-\frac{1}{2}\left(\frac{t-\Delta t}{\sigma\left(\tau_{2}\right)}\right)^{2}\right]
\]
These two pulses defines a space with 8 dimensions to be investigated:
$E_{0}^{(1)}$,  $E_{0}^{(2)}$, $\omega_{1}$, $\omega_{2}$, $CEP_{1}$,
$CEP_{2}$ and $\Delta t$, but also the orientation $P_{1}$ and
$P_{2}$ of the electric field in real space (polarization) parallel
to surface ($s$ or $p$) state for each pulse. 


\paragraph{Effect of the second wavelength}

In particular, it is predicted that effect of the relative $CEP$
for the interfered beams will strongly influence the bicolor transition
probability. 
\begin{figure}
\begin{centering}
\subfloat[]{\begin{centering}
\includegraphics[width=4cm]{\lyxdot \lyxdot /05-QuantumLaP-BicolorprocessingModel/BicolorIrradiation1}
\par\end{centering}

}\subfloat[]{\begin{centering}
\includegraphics[width=4cm]{/media/thibault/PHD/Travail/Prague/Reports/05-QuantumLaP-BicolorprocessingModel/BicolorProcessing-02}
\par\end{centering}

}
\par\end{centering}

\caption{\label{fig:BicolorField}Electric field of the bicolor irradiation.
(a) Source beams. (b) Mixed beams (800 nm and 400 nm). }
\end{figure}
Fig. \ref{fig:BicolorField} shows how two beams of first and second
harmonic are mixing. It is rather clear that a non-periodic modulation
of the electric field will have dramatic consequences on the electron
dynamics, hence changing the absorption probability, in particular
the nonlinear absorption described by Keldysh \cite{Keldysh1965}. 


\paragraph{Effect of the pulse phase (CEP)}

We can also consider two supplementary parameters with adding the
pulse durations $\tau_{1}$ and $\tau_{2}$. These later parameters
are expected to average or increase the importance of $CEP$ parameters.
Fig. \ref{fig:BicolorFieldPhase} shows the effect of changing the
phase of the second pulse. 
\begin{figure}
\begin{centering}
\subfloat[]{\begin{centering}
\includegraphics[width=4cm]{/home/thibault/Documents/octopus/Pulses/800\lyxdot 0nm-Phase0\lyxdot 0+400\lyxdot 0nm-Phase0\lyxdot 0-Delay0\lyxdot 0ps}
\par\end{centering}

}\subfloat[ ]{\begin{centering}
\includegraphics[width=4cm]{/home/thibault/Documents/octopus/Pulses/800\lyxdot 0nm-Phase0\lyxdot 0+400\lyxdot 0nm-Phase36\lyxdot 0-Delay0\lyxdot 0ps}
\par\end{centering}

}
\par\end{centering}

\begin{centering}
\subfloat[]{\begin{centering}
\includegraphics[width=4cm]{/home/thibault/Documents/octopus/Pulses/800\lyxdot 0nm-Phase0\lyxdot 0+400\lyxdot 0nm-Phase72\lyxdot 0-Delay0\lyxdot 0ps}
\par\end{centering}

}\subfloat[]{\begin{centering}
\includegraphics[width=4cm]{/home/thibault/Documents/octopus/Pulses/800\lyxdot 0nm-Phase0\lyxdot 0+400\lyxdot 0nm-Phase108\lyxdot 0-Delay0\lyxdot 0ps}
\par\end{centering}

}
\par\end{centering}

\begin{centering}
\subfloat[]{\begin{centering}
\includegraphics[width=4cm]{/home/thibault/Documents/octopus/Pulses/800\lyxdot 0nm-Phase0\lyxdot 0+400\lyxdot 0nm-Phase144\lyxdot 0-Delay0\lyxdot 0ps}
\par\end{centering}

}\subfloat[]{\begin{centering}
\includegraphics[width=4cm]{/home/thibault/Documents/octopus/Pulses/800\lyxdot 0nm-Phase0\lyxdot 0+400\lyxdot 0nm-Phase180\lyxdot 0-Delay0\lyxdot 0ps}
\par\end{centering}

}
\par\end{centering}

\caption{\label{fig:BicolorFieldPhase}Electric field of the bicolor irradiation.
Mixed beams (800 nm and 400 nm) of various respective phases}
\end{figure}
 Pulse becomes asymmetric in respect of field envelope, which then
could be considered as harmonic in the non-linear regime, i.e., absorption
probability is expected to be also phase depedent. 


\paragraph{Effect of the pulse delay $\Delta t$}

Finally, chirped pulses should be considered, but then may require
another formulation. 
\begin{figure}
\begin{centering}
\subfloat[]{\begin{centering}
\includegraphics[width=4cm]{/home/thibault/Documents/octopus/Pulses/800\lyxdot 0nm-Phase0\lyxdot 0+400\lyxdot 0nm-Phase0\lyxdot 0-Delay-0\lyxdot 02ps}
\par\end{centering}

}\subfloat[ ]{\begin{centering}
\includegraphics[width=4cm]{/home/thibault/Documents/octopus/Pulses/800\lyxdot 0nm-Phase0\lyxdot 0+400\lyxdot 0nm-Phase0\lyxdot 0-Delay-0\lyxdot 012ps}
\par\end{centering}

}
\par\end{centering}

\begin{centering}
\subfloat[]{\begin{centering}
\includegraphics[width=4cm]{/home/thibault/Documents/octopus/Pulses/800\lyxdot 0nm-Phase0\lyxdot 0+400\lyxdot 0nm-Phase0\lyxdot 0-Delay-0\lyxdot 004ps}
\par\end{centering}

}\subfloat[]{\begin{centering}
\includegraphics[width=4cm]{/home/thibault/Documents/octopus/Pulses/800\lyxdot 0nm-Phase0\lyxdot 0+400\lyxdot 0nm-Phase0\lyxdot 0-Delay0\lyxdot 004ps}
\par\end{centering}

}
\par\end{centering}

\begin{centering}
\subfloat[]{\begin{centering}
\includegraphics[width=4cm]{/home/thibault/Documents/octopus/Pulses/800\lyxdot 0nm-Phase0\lyxdot 0+400\lyxdot 0nm-Phase0\lyxdot 0-Delay0\lyxdot 012ps}
\par\end{centering}

}\subfloat[]{\begin{centering}
\includegraphics[width=4cm]{/home/thibault/Documents/octopus/Pulses/800\lyxdot 0nm-Phase0\lyxdot 0+400\lyxdot 0nm-Phase0\lyxdot 0-Delay0\lyxdot 02ps}
\par\end{centering}

}
\par\end{centering}

\caption{\label{fig:BicolorFieldDelay}Electric field of the bicolor irradiation.
Mixed beams (800 nm and 400 nm) of various respective phases}
\end{figure}
 Fig. \ref{fig:BicolorFieldDelay} shows the effect of the delay between
pulses. 


\paragraph{Chirped pulses? }


\subsection{Preparation of the sample}

We generate a silicon slab with using periodic boundary conditions
along the surface, and absorbing boundary conditions in the direction
normal to the surface. 


\paragraph{Which orientation to choose ?}

In the context of nanostructure formation, the effects of crystal
orientation were only slightly investigated. \cite{Crawford2010,Museur,Sedao2014}
It was shown that crystal orientation should not play a role in the
periodic nanostructures of electrodynamic origin. However, it may
play a role in the formation thermodynamically formed periodic structures.
As Si has a diamond-cubic crystalline structure, we choose $\mbox{Si}(100)$,
and Si$(111)$. 


\paragraph{2d-periodic Si slab}

In view of studying the effect of the surface, we build a 2D-periodic
Si lattice with periodic boundary conditions in directions X and Y.
\begin{figure}
\begin{centering}
\subfloat[]{\begin{centering}
\includegraphics[clip,width=4cm]{/home/thibault/Documents/octopus/FHI-aims/Si100-Tier2-Multi/Si-2x1x1-L30/geometry}
\par\end{centering}

}\subfloat[]{\begin{centering}
\includegraphics[clip,width=4cm]{/home/thibault/Documents/octopus/FHI-aims/Si100-Tier2-Multi/Si-4x1x1-L30/geometry}
\par\end{centering}

}
\par\end{centering}

\begin{centering}
\subfloat[]{\begin{centering}
\includegraphics[clip,width=4cm]{/home/thibault/Documents/octopus/FHI-aims/Si100-H-Tier1-Multi/Si-7x2x2-L50/geometry}
\par\end{centering}

}\subfloat[]{\begin{centering}
\includegraphics[clip,width=4cm]{/home/thibault/Documents/octopus/FHI-aims/Si100-H-Tier1-Multi/Si-7x5x5-L50/geometry}
\par\end{centering}

}
\par\end{centering}

\caption{\label{fig:AtomicStructure} Atomic representation of the modeled
slab samples. (a) 2 layers with 2 atoms of Si each. (b) 4 layers composed
of 2 Si atoms each. The number of atomic layers is increased until
convergence of the total energy per atom. (c) Di-hydridized Si(100)
noted $Si(100):H_{2}$. (d) Same as (c) with more atoms. }
\end{figure}
 Fig. \ref{fig:AtomicStructure} shows a primary cell of the Si slab.
To avoid inducing metallic behavior at the interface, H terminations
were added at the surface. The natural oxidation of Si is here neglected.
As samples are considered at ambient temperature, the most probable
surface reconstruction is the dihydridic $1\times1$ occupying the
pending bonds of Si \cite{Mendoza2006}. Fig. 1(c-d) show the Si sample
with such H terminations. 

As we would like to observe the evolution of the wavefunction density
$\left|\psi\right|^{2}$ in the sample, we generate two types of samples.
Slab 1 is the smallest sample where a large grid of k-points allows
for a good convergence of electronic level calculation. Slab 2 is
a second slab describing the same material, where the number of atoms
is increased in view of maximizing the spatially resolved information.
Such procedure does not affect the integrity of calculation, as the
whole physical information is contained into the Kohn-Sham wavefunction
densities $\left|\psi\right|^{2}$. The size of Slab 2 was increased
until exhibiting the same energy per atom as Slab 1. As soon as the
sample was thick enough, overlapping between surface induced energy
levels and bulk energy levels was suppressed, leading to convergence
of the sample to the same total energy. 


\subsubsection{Procedure for building a minimal slab (SLAB 1)}
\begin{description}
\item [{1.}] Calculate the energy per atom, via making a bulk calculation.
$E_{at}=\frac{E_{tot}}{N_{at}}$
\item [{2.}] Increase slab depth until convergence of surface energy (see
Eq. (2) in paper Beschtedt ). Means that both surface are not interacting
(wavefunctions are not overlapping any longer). 
\item [{3.}] Once this slab (minimal number of atoms in the plane) is built,
we start calculating the k-points to calculate band structure. Find
the required number of k points so that electronic total energy per
atom does not change. 
\item [{4.}] Once it is obtained, we found the $N_{e}$ of the ground state
($T=0$ K). The whole observables are available from this knowledge
(DFT).
\end{description}
Fig. \ref{fig:Convergence1} presents the result of convergence for
a vacuum size $L=3$ nm. 
\begin{figure}
\begin{centering}
\subfloat[]{\begin{centering}
\includegraphics[clip,width=8cm]{/home/thibault/Documents/octopus/FHI-aims/Si100-Tier2-Multi/Convergence-Global}
\par\end{centering}

}
\par\end{centering}

\begin{centering}
\subfloat[]{\begin{centering}
\includegraphics[clip,width=8cm]{/home/thibault/Documents/octopus/FHI-aims/Si100-Tier2-Multi/NumberOfLayers}
\par\end{centering}

}
\par\end{centering}

\caption{\label{fig:Convergence1} (a) Convergence of total energy as function
of Si layers and number of k-points. Each layer contains equivalent
of 2 atoms. The number of atomic layers is increased until convergence
of the total energy per atom. }
\end{figure}
 It appears that a k-grid $1\times8\times8$ is sufficient to reach
convergence of the total energy per atom. However, it is unclear if
the number of layers is sufficient. ``Tier1'' and ``Tier2'' are
increasing levels of precision of the function basis of the Si orbitals.
As we would like to improve performances to the maximum, we shall
reduce the precision to get the biggest box possible. 


\subsubsection{Building a large SLAB 2}
\begin{description}
\item [{1.}] From SLAB1 > Step 2, we have a column which should be repeated
in X and Y, using only 1 k-point. 
\item [{2.}] We add atoms in X and Y directions (same number in both directions)
until total energy is converged, and equal to SLAB 1 > Step 4 energy. 
\item [{3.}] Then we know that we have built the proper ground state for
the supercell of Si. 
\end{description}

\section{Simulation results}

We are basically interested in accessing to the temporal density evolution
of electrons \emph{in the state $\boldsymbol{k}$} 
\[
\frac{\partial N_{\boldsymbol{k}}}{\partial t}
\]
promoted from the valence bands to the conduction band of Si as described
by 
\[
N_{\boldsymbol{k}}(t)=\int\psi_{k}^{\dagger}(\boldsymbol{r},t)\,\psi_{k}(\boldsymbol{r},t)\, d\boldsymbol{r}.
\]
As shown recently {[}REF quantum simulation of linear excitation{]},
the linear processes are exciting electrons towards many wave-vectors
of the conduction band. The difficulty is to restrict the number of
carriers to the conduction band of Si, which may also evolve in time.
In particular, at high laser intensity, ps-melting \cite{Siwick2003}
and non-thermal melting \cite{Rousse2001,Sundaram2002} can occur,
where the band gap undergoes a collapsing effect. 

HOW TO SOLVE THIS ? 


\subsection{Band structure of Si slab}

Show the change of band structure as function of number of layers. 


\subsection{Excitation of free-carriers to conduction band with single-color
laser irradiation}


\subsection{Transition probability of electrons to conduction band upon bi-color
laser irradiation}


\subsection{Quantum description of the band gap collapse with increasing laser
fluence}


\section{Conclusion}


\section*{Acknowledgements}

\bibliographystyle{/media/thibault/PHD/Travail/Berlin/Publications/DoublePulseSi/osajnl}
\bibliography{/media/thibault/PHD/Travail/Bibliography/bibliographie_lue}

\end{document}
